% ========== 绪论重点：信号的能量与功率分类（无错版） ==========
\chapter{绪论重点}

\section{信号的能量与功率分类（核心考点）}

\noindent 以下定义基于 \textbf{1$\Omega$ 电阻} 负载，瞬时功率 $P = |f(t)|^2$。

\subsection{1. 信号能量定义}
连续时间信号总能量：
\[
E = \lim_{T \to \infty} \int_{-T}^{T} |f(t)|^2 dt
\]
说明：$E$ 为信号总能量，单位 $\mathrm{J}$（焦耳）。

\subsection{2. 信号平均功率定义}
连续时间信号平均功率：
\[
P = \lim_{T \to \infty} \frac{1}{2T} \int_{-T}^{T} |f(t)|^2 dt
\]
说明：$P$ 为平均功率，单位 $\mathrm{W}$（瓦特）。

\subsection{3. 信号分类规则（正确教材定义）}
根据能量 $E$ 和功率 $P$ 的取值，信号分为三类：
\begin{itemize}
    \item \textbf{能量信号} \\
    判定：$E < \infty$（能量有限），此时必有 $P = 0$。\\
    示例：矩形脉冲、衰减指数信号。
    
    \item \textbf{功率信号} \\
    判定：$P < \infty$（功率有限），此时必有 $E = \infty$。\\
    示例：正弦信号、周期方波、单位阶跃信号。
    
    \item \textbf{非能量非功率信号} \\
    判定：\textbf{既不是能量信号，也不是功率信号} \\
    即：$E \to \infty$ 且 $P \to \infty$。\\
    示例：无界增长信号 $f(t)=t$、增长指数信号 $f(t)=e^{at}(a>0)$。
\end{itemize}

\begin{tcolorbox}[colback=gray!5!white,colframe=black!70!blue,title=考试必背结论]
\textbf{能量信号：} $E<\infty,\ \ P=0$ \\
\textbf{功率信号：} $P<\infty,\ \ E\to\infty$ \\
\textbf{非能量非功率：} 既不是能量信号，也不是功率信号
\end{tcolorbox}